\documentclass[12pt, a4paper]{article}

\usepackage{fontspec}
\usepackage{polyglossia}
\setmainlanguage{russian}
\setotherlanguage{english}
\setmainfont{FreeSerif}
\setsansfont{FreeSans}
\setmonofont{FreeMono}

\usepackage{amsmath, amsthm, amssymb}
\usepackage{mathtools}
\usepackage{enumitem}
\usepackage{microtype}
\usepackage{geometry}
\usepackage{hyperref}
\usepackage{xcolor}
\usepackage{tikz}
\usetikzlibrary{arrows.meta, positioning}

% Паттерны переноса для русского языка недоступны в этом окружении offline
% (hyph-ru отсутствует в установке TeX Live, сеть отключена). Поэтому вместо
% переноса слов разрешаем абзацному алгоритму растягивать межсловные пробелы
% и слегка нарушать «жёсткость» строки — это предотвращает выход текста
% за поля страницы (Overfull \hbox) ценой чуть менее плотного набора.
\tolerance=2000
\emergencystretch=3em
\hbadness=10000
\hfuzz=2pt
\sloppy

\geometry{
  a4paper,
  top=2.5cm, bottom=2.5cm,
  left=3cm, right=3cm
}

\hypersetup{
  colorlinks=true,
  linkcolor=blue!60!black,
  urlcolor=blue!60!black
}

% ---------- theorem environments ----------
\theoremstyle{plain}
\newtheorem{theorem}{Теорема}
\newtheorem{lemma}[theorem]{Лемма}
\newtheorem{corollary}[theorem]{Следствие}

\theoremstyle{definition}
\newtheorem{definition}[theorem]{Определение}
\newtheorem{remark}[theorem]{Замечание}

% ---------- notation shortcuts ----------
\newcommand{\proves}{\vdash}
\renewcommand{\models}{\vDash}
\newcommand{\Lang}{\mathcal{L}}
\newcommand{\Lw}{\mathcal{L}_\omega}
\newcommand{\Struct}{\mathfrak{M}}
\newcommand{\tuple}[1]{\langle #1 \rangle}
\newcommand{\N}{\mathbb{N}}

% ---------- document ----------
\begin{document}

\title{\textbf{Теорема Гёделя о полноте\\
исчисления предикатов первого порядка}}
\author{}
\date{}
\maketitle
\tableofcontents
\bigskip

% ==============================================
\section{Постановка задачи}
% ==============================================

Фиксируем язык первого порядка $\Lang$ с \emph{конечной сигнатурой}:
конечным набором предикатных символов $P_1,\ldots,P_k$,
функциональных символов $f_1,\ldots,f_m$ и констант $a_1,\ldots,a_n$.
Через $\Lang^0$ обозначим множество всех \emph{замкнутых} формул (предложений) языка~$\Lang$.

\begin{theorem}[Гёдель, 1930]\label{thm:completeness}
Пусть $\Gamma \subseteq \Lang^0$ и $\varphi \in \Lang^0$. Тогда
\[
  \Gamma \models \varphi \iff \Gamma \proves \varphi.
\]
\end{theorem}

Направление $(\Leftarrow)$ — \emph{корректность} — доказывается стандартной индукцией
по выводу и мы его опускаем. Основное содержание теоремы~--- направление
$(\Rightarrow)$, \emph{полнота}, которому посвящён этот текст.

Контрапозитивно, полнота эквивалентна следующему утверждению:

\begin{lemma}[о существовании модели]\label{lem:model-existence}
Всякое непротиворечивое множество $\Gamma \subseteq \Lang^0$ выполнимо,
т.\,е.\ существует $\Lang$-структура $\Struct$ такая, что $\Struct \models \varphi$
для всех $\varphi \in \Gamma$.
\end{lemma}

Доказательству этой леммы и посвящён весь дальнейший текст.

% ==============================================
\section{Предварительные сведения}
% ==============================================

\subsection{Структуры и истинность}

\begin{definition}
\emph{$\Lang$-структура} $\Struct = \tuple{M, \{P_i^{\Struct}\}, \{f_j^{\Struct}\}, \{a_l^{\Struct}\}}$
состоит из:
\begin{itemize}
  \item непустого множества $M$ — \emph{носителя};
  \item интерпретации каждого предикатного символа $P_i$ как отношения
        $P_i^{\Struct} \subseteq M^{\text{ar}(P_i)}$;
  \item интерпретации каждого функционального символа $f_j$ как функции
        $f_j^{\Struct}\colon M^{\text{ar}(f_j)} \to M$;
  \item интерпретации каждой константы $a_l$ как элемента $a_l^{\Struct} \in M$.
\end{itemize}
\end{definition}

Истинность замкнутой формулы $\varphi$ в структуре $\Struct$
(обозначение: $\Struct \models \varphi$) определяется стандартной
рекурсией по построению~$\varphi$:
\begin{itemize}
  \item $\Struct \models P_i(t_1,\ldots,t_k)$ если
        $\tuple{t_1^{\Struct},\ldots,t_k^{\Struct}} \in P_i^{\Struct}$;
  \item $\Struct \models \neg\psi$ если $\Struct \not\models \psi$;
  \item $\Struct \models \psi_1 \wedge \psi_2$ если
        $\Struct \models \psi_1$ и $\Struct \models \psi_2$;
  \item $\Struct \models \exists x\,\psi(x)$ если существует $a \in M$ такой,
        что $\Struct \models \psi(a)$ (где $a$ рассматривается как новая константа).
\end{itemize}

\subsection{Непротиворечивость}

Множество $\Gamma$ называется \emph{непротиворечивым}, если $\Gamma \not\proves \bot$,
то есть из него нельзя вывести противоречие.

% ==============================================
\section{Расширение языка: константы Хенкина}
% ==============================================

\subsection{Зачем нужны свидетели}

Пусть формула $\exists x\,\psi(x)$ истинна в модели.
Тогда в модели должен найтись \emph{конкретный элемент},
реализующий~$\psi$. При синтаксическом построении модели такого элемента нет~---
мы строим носитель из термов языка~--- поэтому нам нужно заранее
«пообещать» имя для каждого подобного элемента.

\subsection{Итеративное расширение языка}

Поскольку добавление новых констант порождает новые замкнутые формулы
(в том числе новые экзистенциальные), расширение языка проводится
\emph{итеративно}.

Строим возрастающую цепочку языков:
\[
  \Lang = \Lang_0 \;\subseteq\; \Lang_1 \;\subseteq\; \Lang_2 \;\subseteq\; \cdots
\]
\textbf{Шаг $n \to n+1$.}
Пусть $\Lang_n$ уже построен. Зафиксируем перечисление
всех экзистенциальных предложений языка $\Lang_n$:
\[
  \exists x\,\psi_0^{(n)}(x),\quad \exists x\,\psi_1^{(n)}(x),\quad \ldots
\]

\begin{remark}[Реестр свидетелей]\label{rem:registry}
Среди формул $\exists x\,\psi_k^{(n)}(x)$ могут встретиться такие, для которых
свидетель уже был назначен на одном из предыдущих шагов $i < n$ (например,
формула $\exists x\,\psi(x)$ дословно повторилась, потому что вошла в язык
$\Lang_n$ через несколько разных путей построения). Чтобы избежать
расточительного дублирования констант для одной и той же формулы, ведём
\emph{реестр} $W \colon \{\text{экз.\ формулы}\} \rightharpoonup \{\text{константы}\}$:
прежде чем вводить новую константу для $\exists x\,\psi_k^{(n)}(x)$, проверяем,
не определён ли уже $W(\exists x\,\psi_k^{(n)}(x))$; если определён~--- используем
существующую константу, если нет~--- вводим свежую и заносим её в $W$.

Это техническое уточнение не влияет на корректность дальнейшей конструкции:
без реестра (то есть с дублированием констант для повторяющихся формул)
доказательство леммы~\ref{lem:henkin-consistent} и леммы об истинности
проходит без изменений, поскольку "лишняя" константа~--- это просто
дополнительный, формально независимый элемент будущего носителя.
Реестр лишь делает модель экономнее. Важно не путать это с
\emph{заменой} уже введённой константы~--- константы $c_{n,k}$ один раз
зафиксированы и в дальнейшем неизменны.
\end{remark}

Для каждого такого предложения $\exists x\,\psi_k^{(n)}(x)$, не имеющего
свидетеля в реестре $W$, введём \emph{свежую} константу $c_{n,k}$,
не встречавшуюся ни в каком $\Lang_i$, $i \le n$. Полагаем
\[
  \Lang_{n+1} \;=\; \Lang_n \cup \bigl\{c_{n,k} : k \in \N\bigr\}.
\]

\textbf{Расширенный язык и расширенное множество аксиом.}
Определим
\[
  \Lw \;=\; \bigcup_{n \in \N} \Lang_n.
\]

\begin{lemma}[замкнутость $\Lw$ относительно расширения Хенкина]\label{lem:omega-closed}
Для каждой замкнутой формулы вида $\exists x\,\psi(x)$ языка $\Lw$ в $\Lw$
уже есть свидетельствующая её константа.
\end{lemma}

\begin{proof}
Формула $\exists x\,\psi(x)$ конечна и потому использует лишь конечное
число символов сигнатуры. Поскольку $\Lw = \bigcup_n \Lang_n$~--- объединение
\emph{возрастающей} цепочки языков, все символы, входящие в формулу,
появляются не позднее некоторого конечного шага; значит сама формула
целиком принадлежит некоторому $\Lang_n$. Но тогда $\exists x\,\psi(x)$
встретилась в перечислении экзистенциальных предложений языка $\Lang_n$
на шаге построения $\Lang_{n+1}$ (см.\ замечание~\ref{rem:registry}),
и для неё была введена константа $c_{n,k}$ (или переиспользована уже
имеющаяся) вместе с аксиомой Хенкина $\exists x\,\psi(x) \to \psi(c_{n,k})$.
Эта аксиома лежит в $H \subseteq \Lw^0$.
\end{proof}

\begin{remark}
Именно эта лемма объясняет, почему $\omega$ шагов расширения достаточно
и трансфинитная итерация не нужна: любая экзистенциальная формула,
\emph{какой бы длинной она ни была}, всё равно конечна, а значит
«застревает» на конечном шаге $n$ и получает свидетеля. Понадобиться
$\omega+1$-й шаг мог бы лишь в ситуации, когда в языке появлялись бы
формулы бесконечной длины~--- но язык первого порядка такого не допускает.
\end{remark}

С каждой парой $(n, k)$ свяжем \emph{аксиому Хенкина}:
\[
  \exists x\,\psi_k^{(n)}(x) \;\to\; \psi_k^{(n)}(c_{n,k}).
\]
Обозначим через $H$ множество всех таких аксиом и положим
\[
  \Gamma^* \;=\; \Gamma \cup H.
\]

\begin{lemma}\label{lem:henkin-consistent}
$\Gamma^*$ непротиворечиво в языке $\Lw$.
\end{lemma}

\begin{proof}
Достаточно показать, что добавление \emph{одной} аксиомы Хенкина
к непротиворечивому множеству сохраняет непротиворечивость; тогда
утверждение для всего $H$ следует индукцией по числу аксиом,
фактически использованных в гипотетическом выводе $\bot$
(вывод конечен, поэтому таких аксиом конечное число, и каждая
содержит свою собственную свежую константу, не встречающуюся
в остальных~--- это гарантировано построением).

\smallskip
\noindent\textbf{Базовый шаг.}
Пусть $\Sigma$~--- непротиворечивое множество формул, и пусть
$H \equiv \exists x\,\psi(x) \to \psi(c)$~--- аксиома Хенкина, где
$c$~--- свежая константа, не входящая ни в $\Sigma$, ни в $\psi$.
Покажем, что $\Sigma \cup \{H\}$ непротиворечиво.

Предположим противное:
\[
  \Sigma \cup \bigl\{\exists x\,\psi(x) \to \psi(c)\bigr\} \;\proves\; \bot.
\]

\textit{Шаг 1: теорема дедукции.}
Вынося аксиому Хенкина из посылок, получаем
\[
  \Sigma \;\proves\; \bigl(\exists x\,\psi(x) \to \psi(c)\bigr) \to \bot,
\]
что по законам логики высказываний (введение отрицания, закон
Дунса Скота) равносильно доказательству отрицания самой импликации:
\[
  \Sigma \;\proves\; \exists x\,\psi(x) \;\wedge\; \neg\psi(c).
\]
Тем самым получены \emph{два независимых} вывода из $\Sigma$:
\[
  \text{(i)}\quad \Sigma \proves \exists x\,\psi(x),
  \qquad\qquad
  \text{(ii)}\quad \Sigma \proves \neg\psi(c).
\]

\textit{Шаг 2: замена константы на переменную.}
Рассмотрим вывод (ii). Поскольку $c$ нигде не входит ни в $\Sigma$,
ни в $\psi$ (кроме явно выписанного места подстановки), её присутствие
в дереве вывода чисто номинально: заменив всюду в этом дереве $c$ на
свежую переменную $y$, не связываемую никаким квантором, получаем
вновь корректное дерево вывода (аксиомы и правила логики применимы
к $y$ ровно так же, как к $c$, поскольку нигде не использовалось,
что это именно константа). Следовательно
\[
  \Sigma \;\proves\; \neg\psi(y).
\]

\textit{Шаг 3: обобщение и противоречие.}
Переменная $y$ не входит свободно ни в одну формулу $\Sigma$ (она
была выбрана свежей), поэтому к выводу $\neg\psi(y)$ законно
применимо правило введения квантора всеобщности:
\[
  \Sigma \;\proves\; \forall y\,\neg\psi(y).
\]
По законам обращения с кванторами $\forall y\,\neg\psi(y)$ логически
эквивалентно $\neg\exists y\,\psi(y)$, а с точностью до переименования
связанной переменной~--- формуле $\neg\exists x\,\psi(x)$. Значит
\[
  \text{(iii)}\quad \Sigma \;\proves\; \neg\exists x\,\psi(x).
\]

\textit{Итог.} Сопоставляя (i) и (iii): $\Sigma$ доказывает одновременно
$\exists x\,\psi(x)$ и $\neg\exists x\,\psi(x)$, то есть $\Sigma \proves \bot$~---
противоречие с непротиворечивостью $\Sigma$. Заметим, что в этом
рассуждении свежая константа $c$ исходной теории $\Sigma$ нигде
больше не появляется: противоречие получено строго внутри $\Sigma$.

\smallskip
\noindent\textbf{От одной аксиомы к произвольному конечному множеству.}
Базовый шаг устраняет ровно одну добавленную аксиому Хенкина. Чтобы
распространить результат на всё $H$, нужно действовать аккуратно: вывод
$\Gamma \cup H \proves \bot$ конечен, поэтому фактически использует лишь
конечный список аксиом $H_1, \ldots, H_m \in H$, но \emph{нельзя}
устранить их все одним применением теоремы дедукции.

\begin{remark}[почему «наивный» путь не работает]
Вынося сразу весь список из посылок, получаем
\[
  \Gamma \;\proves\; \neg(H_1 \wedge \cdots \wedge H_m)
  \;\equiv\;
  \Gamma \;\proves\; \neg H_1 \vee \cdots \vee \neg H_m.
\]
Из доказуемости дизъюнкции в классической логике \emph{не следует}
доказуемость хотя бы одного дизъюнкта по отдельности~--- в отличие от
семантического «истинности дизъюнкции», на уровне формального
вывода это не свойство $\proves$. Поэтому распараллелить рассуждение
по ветвям $\Gamma \proves \neg H_i$ нельзя: придётся устранять аксиомы
\emph{последовательно}, а не одним махом.
\end{remark}

Последовательное устранение возможно благодаря следующему наблюдению:
список $H_1, \ldots, H_m$ всегда можно \emph{упорядочить} так, чтобы
константа аксиомы $H_m$ была свежей не только для $\Gamma$, но и для
всех остальных аксиом $H_1, \ldots, H_{m-1}$. Это следует прямо из
построения: константы вводятся по уровням $n = 0, 1, 2, \ldots$, и
внутри одного уровня~--- по разным экзистенциальным формулам, причём
каждая константа $c_{n,k}$ используется только в \emph{своей} аксиоме
Хенкина и никогда не входит в формулы, на основе которых вводились
более ранние константы. Достаточно упорядочить $H_1,\ldots,H_m$ по
неубыванию уровня их констант (а внутри уровня~--- произвольно).

Положим $\Sigma_0 = \Gamma$ и $\Sigma_j = \Gamma \cup \{H_1,\ldots,H_j\}$
для $j = 1,\ldots,m$. По построению константа аксиомы $H_j$ не входит
в $\Sigma_{j-1}$. Применяя базовый шаг к паре $(\Sigma_{j-1}, H_j)$,
получаем: если $\Sigma_{j-1} \cup \{H_j\} = \Sigma_j$ противоречиво, то
противоречиво и $\Sigma_{j-1}$. Применяя это рассуждение последовательно
для $j = m, m-1, \ldots, 1$ («снимая» аксиомы с конца списка),
заключаем: если $\Sigma_m = \Gamma \cup \{H_1,\ldots,H_m\}$
противоречиво, то противоречиво и $\Sigma_0 = \Gamma$.

\smallskip
Таким образом, если $\Gamma \cup H \proves \bot$, то $\Gamma \proves \bot$~---
вопреки предположению о непротиворечивости $\Gamma$. Значит
$\Gamma \cup H = \Gamma^*$ непротиворечиво. Аксиома Хенкина физически
может участвовать в дереве вывода противоречия, но, как показывает
базовый шаг, применённый по одной аксиоме за раз, она структурно
всегда устранима: противоречие, в котором она использована,
сворачивается в более короткое противоречие, в конце концов получаемое
из $\Gamma$ напрямую.
\end{proof}

% ==============================================
\section{Лемма Линденбаума: пополнение до полного множества}
% ==============================================

\begin{definition}
Множество предложений $\Delta \subseteq \Lw^0$ называется
\emph{полным непротиворечивым} (ПН-множеством), если:
\begin{enumerate}[label=(\roman*)]
  \item $\Delta$ непротиворечиво;
  \item для каждого $\varphi \in \Lw^0$ ровно одно из $\varphi,\,\neg\varphi$
        принадлежит $\Delta$.
\end{enumerate}
\end{definition}

\begin{lemma}[Линденбаум]\label{lem:lindenbaum}
Каждое непротиворечивое $\Gamma^* \subseteq \Lw^0$ можно расширить
до ПН-множества $\Delta \supseteq \Gamma^*$.
\end{lemma}

\begin{proof}
Язык $\Lw$ счётен (объединение счётной цепочки счётных языков),
поэтому все предложения $\Lw^0$ можно перечислить:
$\varphi_0, \varphi_1, \varphi_2, \ldots$

Строим возрастающую цепочку:
\[
  \Gamma_0 = \Gamma^*, \qquad
  \Gamma_{n+1} =
  \begin{cases}
    \Gamma_n \cup \{\varphi_n\} & \text{если } \Gamma_n \cup \{\varphi_n\}
    \text{ непротиворечиво,}\\
    \Gamma_n & \text{иначе.}
  \end{cases}
\]

\textbf{Каждый шаг корректен.}
По закону исключённого третьего либо $\Gamma_n \cup \{\varphi_n\}$
непротиворечиво, либо нет~--- третьего не дано.
(Заметим: здесь мы не требуем \emph{разрешимости} проверки непротиворечивости~---
достаточно её \emph{существования} как дихотомии.)

Полагаем $\Delta = \bigcup_{n} \Gamma_n$.

\textbf{$\Delta$ непротиворечиво.}
Пусть $\Delta \proves \bot$. Тогда этот вывод конечен и использует
лишь формулы из некоторого $\Gamma_N$~--- противоречие.

\textbf{$\Delta$ полно.}
Пусть $\varphi = \varphi_n$. Если $\Gamma_n \cup \{\varphi_n\}$ непротиворечиво,
то $\varphi_n \in \Gamma_{n+1} \subseteq \Delta$.
Иначе $\Gamma_n \cup \{\varphi_n\} \proves \bot$, значит $\Gamma_n \proves \neg\varphi_n$,
откуда $\neg\varphi_n \in \Delta$.
Таким образом $\Delta$ содержит ровно одно из $\varphi_n, \neg\varphi_n$.
\end{proof}

\begin{remark}
В счётном случае конструкция полностью явная (никакого трансфинита
и никакой аксиомы выбора). Для несчётных языков потребовалась бы
лемма Цорна, но мы её не используем.
\end{remark}

\begin{remark}[Неконструктивность $\Delta$ и множественность моделей]
Объект $\Delta$ существует в силу классических принципов (ЗИТ и счётной
итерации), но в общем случае \emph{неэффективен}: мы не можем алгоритмически
определить, какие именно формулы в него вошли, если только они не выводимы
из $\Gamma^*$ явно. На каждом шаге мы делаем неконструктивный выбор между
$\varphi_n$ и $\neg\varphi_n$ — существование такого выбора гарантировано
ЗИТ, но конкретный исход нам недоступен.

Следствием этого является то, что разные «достройки» $\Gamma$ до
ПН-множества дают в общем случае \emph{неизоморфные} модели. Именно поэтому
непротиворечивая теория как правило имеет много различных моделей.
Это не дефект конструкции, а отражение реальности: теорема Гёделя
о полноте доказывает \emph{существование} модели, а не её
единственность или конструктивную предъявимость.
\end{remark}

% ==============================================
\section{Построение термовой модели}
% ==============================================

Далее $\Delta$~--- фиксированное ПН-множество, содержащее $\Gamma^*$.

\subsection{Носитель}

Пусть $T$ — множество всех \emph{замкнутых термов} языка $\Lw$
(термов без свободных переменных). Вводим на $T$ отношение:
\[
  s \sim t \;\stackrel{\text{def}}{\iff}\; (s = t) \in \Delta.
\]

\begin{lemma}
$\sim$ является отношением эквивалентности.
\end{lemma}

\begin{proof}
Рефлексивность: $(t = t) \in \Delta$, так как $\proves t = t$.
Симметричность и транзитивность аналогично следуют из
того, что $\Delta$ содержит все аксиомы равенства.
\end{proof}

Носитель модели:
\[
  M \;=\; T/{\sim}, \qquad [t] = \{s \in T : s \sim t\}.
\]

\subsection{Природа термовой модели}

Ключевое наблюдение: носитель $M = T/{\sim}$ состоит из \emph{термов},
а не из каких-то абстрактных объектов. Поэтому интерпретация символов
языка — это не «вычисление в заранее данной структуре», а
\emph{определение по синтаксису}.

Иными словами: мы не знаем и не обязаны знать, «что такое» константа
Хенкина $c_{n,k}$ как объект. Она \emph{есть} элемент модели $[c_{n,k}]$
по определению. Вопрос «выполнено ли $P(c_{n,k})$?» имеет ответ:
«да, если $P(c_{n,k}) \in \Delta$; нет, если $\neg P(c_{n,k}) \in \Delta$»~---
и ровно один из этих вариантов реализован, поскольку $\Delta$ полно.
Никакого дополнительного «знания» о $c_{n,k}$ не требуется и не предполагается.

Аналогично, значение функции $f^{\Struct}([t_1],\ldots,[t_r])$ — это просто
класс терма $[f(t_1,\ldots,t_r)]$, то есть применение функционального символа
на уровне синтаксиса термов. Семантика полностью редуцирована к синтаксису
через принадлежность $\Delta$.

\subsection{Интерпретация символов}

Прежде чем задавать интерпретацию, нужно зафиксировать, что лежит
в основании формальной системы: мы предполагаем, что исчисление
предикатов с равенством включает (помимо рефлексивности,
симметричности и транзитивности) \emph{схему аксиом конгруэнтности}
для каждого функционального символа $f_j$ арности $r$ и каждого
предикатного символа $P_i$ арности $r$:
\[
  \forall \vec x\, \forall \vec y\;
  \Bigl(\textstyle\bigwedge_{i=1}^{r} x_i = y_i \;\to\; f_j(\vec x) = f_j(\vec y)\Bigr),
\]
\[
  \forall \vec x\, \forall \vec y\;
  \Bigl(\textstyle\bigwedge_{i=1}^{r} x_i = y_i \;\to\; (P_i(\vec x) \leftrightarrow P_i(\vec y))\Bigr).
\]
Это стандартное предположение (Эндертон, Mendelson и др.): без него
равенство было бы лишь произвольным отношением эквивалентности,
не согласованным с остальной сигнатурой, и определения ниже были бы
некорректны. Поскольку эти аксиомы логические, они выводимы из
пустого множества посылок, а значит входят в $\Delta$ (как и в любое
полное непротиворечивое множество, содержащее все теоремы исчисления).

\begin{itemize}
  \item \textbf{Константы.} $a_l^{\Struct} = [a_l]$ для каждой константы
        $a_l \in \Lang$ (а также для констант Хенкина $c_{n,k}$).

  \item \textbf{Функциональные символы.}
        \[
          f_j^{\Struct}([t_1], \ldots, [t_r]) \;:=\; [f_j(t_1, \ldots, t_r)].
        \]

  \item \textbf{Предикатные символы.}
        \[
          P_i^{\Struct}([t_1], \ldots, [t_r]) = \top
          \;\iff\; P_i(t_1,\ldots,t_r) \in \Delta.
        \]
\end{itemize}

\begin{lemma}[корректность определений]\label{lem:welldefined}
Оба определения не зависят от выбора представителей классов
эквивалентности.
\end{lemma}

\begin{proof}
Пусть $s_i \sim t_i$ для всех $i = 1,\ldots,r$, то есть
$(s_i = t_i) \in \Delta$.

\textbf{Функции.} Аксиома конгруэнтности для $f_j$ есть логическая
теорема, значит лежит в $\Delta$. Применяя её (через схему общности
и modus ponens, используя $(s_i = t_i) \in \Delta$ для каждого $i$),
получаем
$\bigl(f_j(s_1,\ldots,s_r) = f_j(t_1,\ldots,t_r)\bigr) \in \Delta$,
то есть $f_j(s_1,\ldots,s_r) \sim f_j(t_1,\ldots,t_r)$, а значит
$[f_j(s_1,\ldots,s_r)] = [f_j(t_1,\ldots,t_r)]$.

\textbf{Предикаты.} Аналогично, из аксиомы конгруэнтности для $P_i$
и $(s_i=t_i)\in\Delta$ получаем
$\bigl(P_i(s_1,\ldots,s_r) \leftrightarrow P_i(t_1,\ldots,t_r)\bigr) \in \Delta$.
Поскольку $\Delta$ полно и замкнуто относительно вывода, отсюда следует
$P_i(s_1,\ldots,s_r) \in \Delta \iff P_i(t_1,\ldots,t_r) \in \Delta$,
то есть значение $P_i^{\Struct}$ действительно не зависит от выбора
представителей.
\end{proof}

\begin{remark}
Это не «теоретико-множественная магия»~--- корректность существенно
использует то, что аксиомы конгруэнтности явно входят в логический
базис исчисления предикатов с равенством. Если бы равенство было
произвольным конгруэнтным отношением, заданным без этой схемы аксиом,
определения $f^{\Struct}$ и $P^{\Struct}$ нуждались бы в отдельном
обосновании или попросту были бы некорректны.
\end{remark}

% ==============================================
\section{Лемма об истинности}
% ==============================================

\begin{lemma}[об истинности]\label{lem:truth}
Для каждого замкнутого $\varphi \in \Lw^0$:
\[
  \Struct \models \varphi \;\iff\; \varphi \in \Delta.
\]
\end{lemma}

\begin{proof}
Индукция по сложности~$\varphi$.

\medskip
\noindent\textbf{Атомный случай $\varphi = P_i(t_1,\ldots,t_r)$.}
По определению $P_i^{\Struct}$:
$\Struct \models P_i(t_1,\ldots,t_r) \iff P_i(t_1,\ldots,t_r) \in \Delta$.

\medskip
\noindent\textbf{Случай $\varphi = \neg\psi$.}
По полноте $\Delta$: $\neg\psi \in \Delta \iff \psi \notin \Delta$.
По индукционному предположению: $\psi \notin \Delta \iff \Struct \not\models \psi$.
Итого: $\neg\psi \in \Delta \iff \Struct \models \neg\psi$. \checkmark

\medskip
\noindent\textbf{Случай $\varphi = \psi_1 \wedge \psi_2$.}
$\psi_1 \wedge \psi_2 \in \Delta$ iff $\psi_1 \in \Delta$ и $\psi_2 \in \Delta$
(так как $\Delta$ замкнуто относительно вывода в обе стороны
и содержит аксиомы для конъюнкции),
что по индукции равносильно $\Struct \models \psi_1 \wedge \psi_2$. \checkmark

\medskip
\noindent\textbf{Ключевой случай $\varphi = \exists x\,\psi(x)$.}

$(\Rightarrow)$ Пусть $\Struct \models \exists x\,\psi(x)$.
Тогда существует $[t] \in M$ такой, что $\Struct \models \psi(t)$.
По индукции $\psi(t) \in \Delta$, откуда $\proves \exists x\,\psi(x)$,
следовательно $\exists x\,\psi(x) \in \Delta$.

$(\Leftarrow)$ Пусть $\exists x\,\psi(x) \in \Delta$.
По построению существует аксиома Хенкина
$\exists x\,\psi(x) \to \psi(c)$, которая лежит в $\Gamma^* \subseteq \Delta$.
Modus ponens: $\psi(c) \in \Delta$.
По индукционному предположению $\Struct \models \psi(c)$,
а значит $\Struct \models \exists x\,\psi(x)$. \checkmark
\end{proof}

% ==============================================
\section{Завершение доказательства}
% ==============================================

Перед тем как завершить доказательство, нам понадобится общий факт о
структурах: истинность формулы зависит только от интерпретации тех
символов, которые в неё реально входят, а не от всей сигнатуры языка,
в котором структура задана.

\begin{lemma}[о совпадении интерпретаций]\label{lem:reduct}
Пусть $\Lang' \subseteq \Lang''$~--- языки (сигнатура $\Lang'$ является
частью сигнатуры $\Lang''$), и пусть
$\Struct'' = \tuple{M, I''}$~--- $\Lang''$-структура, где $I''$~---
её функция интерпретации символов. Определим
\emph{редукт} (обеднение) $\Struct''$ до сигнатуры $\Lang'$ как
\[
  \Struct''|_{\Lang'} \;=\; \tuple{M, I'},
  \qquad \text{где } I' = I''\big|_{\text{символы } \Lang'}
\]
(тот же носитель $M$, а интерпретация $I'$~--- это сужение $I''$ только
на символы $\Lang'$; интерпретации символов из $\Lang'' \setminus \Lang'$
просто отбрасываются).

Тогда для любой формулы $\varphi$ языка $\Lang'$ (то есть не содержащей
символов вне $\Lang'$):
\[
  \Struct'' \models \varphi
  \quad\iff\quad
  \Struct''|_{\Lang'} \models \varphi.
\]
\end{lemma}

\begin{proof}
Индукция по построению формулы $\varphi$. На атомном шаге $\varphi$ имеет
вид $P(t_1,\ldots,t_r)$ или $s=t$ для символов $P$ и термов из $\Lang'$;
истинность вычисляется через $I''(P)$ либо через $I'(P) = I''(P)$~---
одно и то же отношение, поскольку $P \in \Lang'$. Аналогично для термов:
значение $t^{\Struct''}$ и $t^{\Struct''|_{\Lang'}}$ совпадают, так как
все функциональные символы и константы, входящие в $t$, лежат в $\Lang'$
и потому интерпретируются одинаково в $I''$ и в $I'$. Шаги индукции для
$\neg, \wedge, \exists$ дословно переносят определение истинности с
$\Struct''$ на $\Struct''|_{\Lang'}$ и обратно, поскольку само
определение истинности на этих связках не упоминает сигнатуру явно.
\end{proof}

\begin{remark}
Содержательно: ограничение сигнатуры~--- чисто структурная операция
отсечения неиспользуемых интерпретаций, не меняющая ни носитель $M$,
ни поведение структуры на формулах, которые эти отброшенные символы
не используют.
\end{remark}

\begin{proof}[Доказательство леммы~\ref{lem:model-existence}]
Пусть $\Gamma \subseteq \Lang^0$ непротиворечиво.

\begin{enumerate}[label=\textbf{Шаг \arabic*.}, leftmargin=*, itemsep=4pt]
  \item \textbf{Расширение языка.}
        Строим цепочку $\Lang_0 \subseteq \Lang_1 \subseteq \cdots$
        и $\Lw = \bigcup_n \Lang_n$ с константами Хенкина,
        формируем $\Gamma^* = \Gamma \cup H$.
        По лемме~\ref{lem:henkin-consistent} $\Gamma^*$ непротиворечиво.

  \item \textbf{Пополнение по Линденбауму.}
        По лемме~\ref{lem:lindenbaum} расширяем $\Gamma^*$
        до ПН-множества $\Delta \subseteq \Lw^0$.

  \item \textbf{Термовая модель.}
        Строим $\Lw$-структуру $\Struct$ с носителем $T/{\sim}$
        и интерпретацией символов, как описано в разделе~5.

  \item \textbf{Применение леммы об истинности.}
        По лемме~\ref{lem:truth} для каждого $\varphi \in \Delta$
        имеем $\Struct \models \varphi$.
        Поскольку $\Gamma \subseteq \Gamma^* \subseteq \Delta$,
        получаем $\Struct \models \Gamma$.

  \item \textbf{Редукт на исходный язык.}
        Применим лемму~\ref{lem:reduct} к языкам $\Lang \subseteq \Lw$ и
        структуре $\Struct$:
        \begin{enumerate}[label=(\alph*)]
          \item исходная теория $\Gamma$ сформулирована в языке $\Lang$;
          \item мы построили расширенный язык $\Lw \supseteq \Lang$,
                добавив константы Хенкина;
          \item термовая структура $\Struct$ задана в языке $\Lw$, и
                по предыдущему шагу $\Struct \models \Gamma$;
          \item все формулы $\Gamma$ записаны в языке $\Lang$ (в них нет
                ни одной константы Хенкина), поэтому возьмём редукт
                $\Struct|_{\Lang}$~--- ту же структуру с носителем $M$,
                но с забытой интерпретацией символов $c_{n,k}$;
          \item по лемме~\ref{lem:reduct}, применённой к каждой формуле
                $\varphi \in \Gamma$, получаем $\Struct|_{\Lang} \models \Gamma$.
        \end{enumerate}
\end{enumerate}
\end{proof}

\begin{corollary}[Теорема Гёделя о полноте]
$\Gamma \models \varphi \implies \Gamma \proves \varphi$.
\end{corollary}

\begin{proof}
Если $\Gamma \not\proves \varphi$, то $\Gamma \cup \{\neg\varphi\}$ непротиворечиво.
По лемме~\ref{lem:model-existence} существует модель $\Struct$ такая,
что $\Struct \models \Gamma$ и $\Struct \models \neg\varphi$,
то есть $\Struct \not\models \varphi$.
Следовательно $\Gamma \not\models \varphi$.
\end{proof}

% ==============================================
\section{Пример: теорема Акса–Гротендика}
% ==============================================

Одно из наиболее эффектных применений теории моделей — доказательство
следующего факта, принадлежащего алгебраической геометрии.

\begin{theorem}[Акс–Гротендик]
Любое инъективное полиномиальное отображение $f \colon \mathbb{C}^n \to \mathbb{C}^n$
является сюръективным.
\end{theorem}

Примечательно, что для самого доказательства не требуется ни топологии,
ни анализа, ни никакого явного построения прообраза~--- только теория моделей.

\begin{proof}[Доказательство (Акс, 1968)]

\textbf{Шаг 1: формализация.}
Утверждение «каждое инъективное полиномиальное отображение $\mathbb{F}^n \to \mathbb{F}^n$
сюръективно» для фиксированной размерности $n$ и фиксированных степеней полиномов
$(d_1,\ldots,d_n)$ выражается одной формулой первого порядка в языке колец:
\[
  \sigma_{n,\vec d} \;\equiv\;
  \forall \vec a\; \Bigl[
    \bigl(\forall \vec x\,\forall \vec y\;
      (f_{\vec a}(\vec x)=f_{\vec a}(\vec y) \to \vec x = \vec y)\bigr)
    \to
    \forall \vec z\;\exists \vec x\;f_{\vec a}(\vec x) = \vec z
  \Bigr],
\]
где $f_{\vec a}$ обозначает отображение с набором коэффициентов $\vec a$.
Полное утверждение теоремы для $\mathbb{C}$ есть $\mathbb{C} \models \sigma_{n,\vec d}$
для всех $n$ и $\vec d$.

\textbf{Шаг 2: случай конечных полей.}
Если $\mathbb{F}$ конечно, то $\mathbb{F}^n$ конечно. Инъекция конечного
множества в себя сюръективна по соображениям мощности. Следовательно,
$\mathbb{F} \models \sigma_{n,\vec d}$ для любого конечного поля $\mathbb{F}$
и любых $n,\vec d$.

\textbf{Шаг 3: теория $\mathrm{ACF}_0$ и признак Вота.}
Напомним, что теория \emph{алгебраически замкнутых полей нулевой характеристики}
($\mathrm{ACF}_0$) описывается аксиомами алгебраически замкнутого поля вместе со
схемой аксиом $\neg(1+\cdots+1)=0$ (по одной на каждое простое $p$). Эта теория
\emph{полна}: по \emph{признаку Вота} теория счётной сигнатуры, не имеющая конечных
моделей и $\kappa$-категоричная для некоторого бесконечного $\kappa$, полна.
$\mathrm{ACF}_0$ категорична в любой несчётной мощности (все алгебраически замкнутые
поля нулевой характеристики одной несчётной мощности изоморфны), откуда полнота.

В частности, $\mathbb{C}$~--- единственная (с точностью до изоморфизма) несчётная
модель $\mathrm{ACF}_0$ мощности $2^{\aleph_0}$, и $\mathrm{ACF}_0 \models \varphi$
равносильно $\mathbb{C} \models \varphi$ для любого предложения $\varphi$.

\textbf{Шаг 4: перенос через теорему компактности.}
Пусть дано конкретное $\sigma_{n,\vec d}$.

Предположим, что $\mathbb{C} \not\models \sigma_{n,\vec d}$.
Поскольку $\mathrm{ACF}_0$ полна и $\mathbb{C} \models \mathrm{ACF}_0$, это
означало бы $\mathrm{ACF}_0 \models \neg\sigma_{n,\vec d}$. По теореме о полноте
(теорема~\ref{thm:completeness}) существует формальный вывод
$\mathrm{ACF}_0 \vdash \neg\sigma_{n,\vec d}$. Этот вывод конечен, поэтому
использует лишь конечное число аксиом $\mathrm{ACF}_0$~--- в частности,
лишь конечное число аксиом $\neg p\cdot 1 = 0$, скажем, для простых
$p \le p_0$.

Значит, $\neg\sigma_{n,\vec d}$ выводима уже из
\[
  \mathrm{ACF}_{>p_0} \;:=\; \mathrm{ACF} \;\cup\;
  \{\neg p \cdot 1 = 0 : p \le p_0\},
\]
и потому $\mathrm{ACF}_{>p_0} \models \neg\sigma_{n,\vec d}$.
Но для любого простого $p > p_0$ алгебраическое замыкание $\overline{\mathbb{F}_p}$
является моделью $\mathrm{ACF}_{>p_0}$, а значит
$\overline{\mathbb{F}_p} \models \neg\sigma_{n,\vec d}$. При этом
$\overline{\mathbb{F}_p}$ содержит конечные подполя $\mathbb{F}_{p^k}$, и
ограничение рассуждения Шага~2 даёт $\overline{\mathbb{F}_p} \models \sigma_{n,\vec d}$~---
противоречие.

Следовательно, $\mathbb{C} \models \sigma_{n,\vec d}$.
\end{proof}

\begin{remark}[методологическое значение]
Доказательство показывает принципиальную черту модельно-теоретического метода:
из полноты $\mathrm{ACF}_0$ и теоремы о полноте заключаем, что существует
чисто алгебраический формальный вывод теоремы Акса–Гротендика из аксиом теории
колец, хотя этот вывод нигде не предъявляется и на практике был бы
чрезвычайно сложен. Таким образом, модельная теория устанавливает \emph{существование}
доказательства, не строя его явно~--- и именно поэтому такой метод называют
\emph{перефразированием} или \emph{переносом}: факт, доступный в одном
классе структур, переносится в другой через семантическое посредство.
\end{remark}

% ==============================================
\section*{Иллюстрация: ход конструкции на сквозном примере}
% ==============================================

Следующая схема собирает воедино три шага конструкции для языка
$\mathcal{L} = \{=,\,{<},\,0\}$ из раздела~7: рост языка по уровням,
переход к полному непротиворечивому множеству и итоговый носитель
термовой модели с уже проявившимся в нём порядком.

\begin{center}
\begin{tikzpicture}[
  scale=0.88, every node/.style={scale=0.88},
  font=\small,
  level/.style={draw, rounded corners, minimum width=2.6cm, minimum height=0.6cm, align=center},
  arrow/.style={-{Stealth[length=2.2mm]}, thick},
  annot/.style={font=\tiny, align=left, text width=2.8cm},
]

% ============ COLUMN 1: language tower ============
\node[level, fill=blue!8]   (L0) at (0,0)   {$\mathcal{L}_0 = \{=,<,0\}$};
\node[level, fill=teal!10]  (L1) at (0,1.6) {$\mathcal{L}_1$};
\node[level, fill=teal!10]  (L2) at (0,3.2) {$\mathcal{L}_2$};
\node[level, fill=teal!25]  (Lw) at (0,4.8) {$\mathcal{L}_\omega = \bigcup_n \mathcal{L}_n$};

\draw[arrow] (L0) -- (L1);
\node[annot, anchor=west] at (1.5,0.8) {$+\,c_{0,0},c_{0,1}$\\свидетели для $\exists x(0{<}x)$};

\draw[arrow] (L1) -- (L2);
\node[annot, anchor=west] at (1.5,2.4) {$+\,c_{1,0},c_{1,1},\ldots$\\свидетели формул $\mathcal{L}_1$};

\draw[arrow, dashed] (L2) -- (Lw);
\node[font=\tiny, anchor=west] at (1.5,4.0) {$\omega$ шагов};

\node[font=\footnotesize\bfseries] at (0,-0.85) {Шаг 1};
\node[font=\footnotesize] at (0,-1.2) {расширение языка};

% ============ COLUMN 2: Lindenbaum ============
\node[level, fill=orange!15, minimum width=1.6cm] (Delta) at (6.0,4.8) {$\Delta$};

\draw[arrow, ultra thick, blue!45!black] (Lw) -- (Delta);
\node[font=\footnotesize] at (4.55,5.35) {Линденбаум};
\node[font=\tiny] at (4.55,5.0) {$\Gamma^* \rightsquigarrow \Delta$};
\node[annot, anchor=north] at (4.45,4.55)
  {полное непротиво-\\речивое множество};

\node[font=\footnotesize\bfseries] at (6.0,-0.85) {Шаг 2};
\node[font=\footnotesize] at (6.0,-1.2) {пополнение Линденбаума};

% ============ COLUMN 3: model carrier ============
\draw[arrow, ultra thick, teal!50!black] (Delta) -- (8.7,4.8);
\node[font=\footnotesize] at (10.0,5.35) {синтез модели};

\begin{scope}[shift={(9.0,2.05)}]
\draw[thick] (0,0) -- (4.0,0);
\draw (0,-0.1)    -- (0,0.1);    \node[font=\tiny, below=2pt] at (0,-0.15)    {$[0]$};
\draw (0.75,-0.1) -- (0.75,0.1); \node[font=\tiny, below=2pt] at (0.75,-0.15) {$[c_{1,2}]$};
\draw (1.5,-0.1)  -- (1.5,0.1);  \node[font=\tiny, below=2pt] at (1.5,-0.15)  {$[c_{0,0}]$};
\draw (2.2,-0.1)  -- (2.2,0.1);  \node[font=\tiny, below=2pt] at (2.2,-0.15)  {$[c_{1,0}]$};
\draw (2.95,-0.1) -- (2.95,0.1); \node[font=\tiny, below=2pt] at (2.95,-0.15) {$[c_{0,1}]$};
\draw (3.7,-0.1)  -- (3.7,0.1);  \node[font=\tiny, below=2pt] at (3.7,-0.15)  {$[c_{1,1}]$};
\node[font=\tiny] at (4.35,0) {$\cdots$};
\node[font=\footnotesize] at (2.0,0.5) {носитель $T/{\sim}$};
\end{scope}

\node[annot, anchor=north, text width=4cm] at (11.0,3.7)
  {$P_<^{\mathfrak M}([s],[t])=\top$\\[2pt] $\iff (s{<}t)\in\Delta$};

\node[font=\footnotesize\bfseries] at (11.0,-0.85) {Шаг 3};
\node[font=\footnotesize] at (11.0,-1.2) {термовая модель};

\end{tikzpicture}
\end{center}

% ==============================================
\section*{Резюме конструкции}
% ==============================================

\begin{center}
\footnotesize
\begin{tabular}{cp{12cm}}
  \hline
  \textbf{Шаг} & \textbf{Что делаем} \\
  \hline
  1 & Расширяем язык: добавляем константы Хенкина $c_{n,k}$ итеративно
      (с реестром свидетелей, замечание~\ref{rem:registry}) \\
  2 & Пополняем $\Gamma$ аксиомами Хенкина~--- получаем $\Gamma^*$,
      непротиворечивое по лемме~\ref{lem:henkin-consistent} \\
  3 & Алгоритмом Линденбаума строим ПН-множество $\Delta \supseteq \Gamma^*$ \\
  4 & Носитель модели~--- классы эквивалентности замкнутых термов $\Lw$ \\
  5 & Предикаты и функции задаём через $\Delta$; корректность даёт
      лемма~\ref{lem:welldefined} (аксиомы конгруэнтности) \\
  6 & Лемма об истинности (индукция по сложности) даёт $\Struct \models \Delta$ \\
  7 & Редукт $\Struct|_{\Lang}$ (лемма~\ref{lem:reduct}) является моделью исходного $\Gamma$ \\
  \hline
\end{tabular}
\end{center}

\vfill
\begin{center}
\rule{0.3\textwidth}{0.4pt}\\[0.5em]
{\footnotesize\color{gray} \copyright\ Nikolai Kazimirov 2026 \quad\textbar\quad \url{https://mathem.at}}
\end{center}

\end{document}
